\section{Exact sixth- and seventh-index saturation thresholds}
\label{higherzeros:sec:thresholds}
The incoming fifth-index method suggests a reproducible extension of the
classical zero classification. Its next two first-derivative counts were
reported only as numerical conjectures. Exact critical-value descent now
settles both indices at every derivative order, not only the sampled first
order. The argument retains the existing global bound and adds finite
rational signs; it does not infer completeness from an approximate scan.

\begin{theorem}[The next two classical transitions]
\label{higherzeros:thm:thresholds}
Let $N_{n,k}$ count the positive zeros of $\gamma_n^{(k)}$, with
multiplicity. For every integer $k\ge1$,
\[
 N_{6,k}=\begin{cases}4,&1\le k\le3,\\6,&k\ge4,\end{cases}
 \qquad
 N_{7,k}=\begin{cases}5,&1\le k\le4,\\7,&k\ge5.\end{cases}
\]
Every zero is simple. Thus the least classical saturation thresholds are
$K_6^{\rm cl}=4$ and $K_7^{\rm cl}=5$. At and above each threshold,
the next-order zeros strictly interlace to the right of the current ones.
\end{theorem}

Retain $U_{n,k}=(-1)^{n+k}\gamma_n^{(k)}$ and the sign-preserving
normalization $V_{n,k}=a^{k+1}U_{n,k}/k!$ of the previous section.
The endpoint expansions give $U_{n,k}>0$ near zero,
$\operatorname{sgn}U_{n,k}=(-1)^n$ near infinity, and $U_{n,k}\to0$
there. Also $U_{n,k}'=-U_{n,k+1}$.

\begin{lemma}[Complete descent from critical-value signs]
\label{higherzeros:lem:descent}
Suppose all positive zeros $c_1<\cdots<c_r$ of $U_{n,k+1}$ are
simple, and $U_{n,k}(c_j)$ has a nonzero certified sign $s_j$ at
each one. Put $s_0=1$. Then every positive zero of $U_{n,k}$ is
simple and
\[
 N_{n,k}=\#\{0\le j<r:s_j\ne s_{j+1}\}.
\]
There is no zero on the final interval $(c_r,\infty)$.
\end{lemma}
\begin{proof}
On each open interval cut out by the complete critical-point list,
$U_{n,k}'=-U_{n,k+1}$ is nonzero and has constant sign. Hence
$U_{n,k}$ is strictly monotone there. Its limits or endpoint values
show one crossing exactly when the stated signs differ. It cannot
vanish at a critical point, since each critical value is nonzero.
On the final interval it is strictly monotone with limit zero; a
nonzero strictly monotone function tending to zero cannot cross zero
there. Every crossing lies where the derivative is nonzero, so it
is simple. These intervals cover the whole positive axis.
\end{proof}

\begin{table}[htbp]
\centering\small
\begin{tabular}{cccrcc}
\toprule
$n$&$k$&$j$&Lower endpoint $\ell$&$V_{n,k}(\ell)$&$V_{n,k}(\ell+10^{-14})$\\
\midrule
6&4&1&0.95661907907155&$+$&$-$\\
6&4&2&1.07179934147775&$-$&$+$\\
6&4&3&1.51367699674930&$+$&$-$\\
6&4&4&3.14848332199473&$-$&$+$\\
6&4&5&16.34597601190147&$+$&$-$\\
6&4&6&1389.12290266995208&$-$&$+$\\
7&5&1&0.97564797285268&$+$&$-$\\
7&5&2&1.03011037496236&$-$&$+$\\
7&5&3&1.39376738084727&$+$&$-$\\
7&5&4&2.31266214479726&$-$&$+$\\
7&5&5&5.93513467415295&$+$&$-$\\
7&5&6&38.92782595180997&$-$&$+$\\
7&5&7&5297.39431175614964&$+$&$-$\\
\bottomrule
\end{tabular}
\caption{All saturated-order root brackets. Each displayed decimal is an
exact rational number, and each bracket has width $10^{-14}$. The signs
are determined by outward integer intervals, rather than floating-point
root approximations.}
\label{higherzeros:tab:terminal}
\end{table}

\begin{table}[htbp]
\centering\small
\begin{tabular}{ccclc}
\toprule
$n$&Known critical order&Predecessor order&Signs at all critical points&Count\\
\midrule
6&4&3&$+,+,-,+,-,+$&4\\
6&3&2&$-,+,-,+$&4\\
6&2&1&$-,+,-,+$&4\\
\midrule
7&5&4&$+,+,-,+,-,+,-$&5\\
7&4&3&$-,+,-,+,-$&5\\
7&3&2&$-,+,-,+,-$&5\\
7&2&1&$-,+,-,+,-$&5\\
\bottomrule
\end{tabular}
\caption{The complete critical-value descent. Each predecessor sign is
certified uniformly over the entire rational root bracket of its derivative.}
\label{higherzeros:tab:descent}
\end{table}

\begin{proof}[Proof of Theorem~\ref{higherzeros:thm:thresholds}]
The terminal brackets in Table~\ref{higherzeros:tab:terminal} have opposite
strict endpoint signs, so give six distinct zeros at $(n,k)=(6,4)$ and
seven at $(7,5)$. The global bound with multiplicity, Theorem~\ref{zeros:thm:globalzeros},
makes these complete lists and all their zeros simple. On every terminal
bracket the preceding derivative has the sign recorded in the first
applicable row of Table~\ref{higherzeros:tab:descent}. Thus
Lemma~\ref{higherzeros:lem:descent} gives exactly four sixth-index roots at
order three and five seventh-index roots at order four.

The same exact certificate file supplies disjoint brackets with opposite
endpoint signs for every one of these predecessor roots. Since the preceding
descent already proves the total count, these brackets form the complete
critical-point lists for the next step. Their whole-interval predecessor
signs are the next rows of Table~\ref{higherzeros:tab:descent}. Iterating the
lemma reaches order one, with the counts displayed in the theorem and
simple zeros throughout. This use of lower-order brackets has no circular
completeness assumption: their completeness is established before the next
descent step.

At a saturated order there are $n$ distinct simple zeros. Rolle's theorem
supplies $n-1$ next-order zeros between adjacent ones and one after the
last, because the function has a nonzero sign there and tends back to
zero. The same global bound makes these exactly $n$ simple zeros, with
strict right interlacing. Induction covers every larger order. The lower
counts prove minimality of the two thresholds.
\end{proof}

\paragraph{The exact evidence and its analytic contract.}
The finite proof data contain 45 rational root brackets: 18 for the sixth
index and 27 for the seventh. They require \textbf{90} endpoint enclosures
and \textbf{36} whole critical-interval enclosures. The first implementation
uses the generic Euler--Maclaurin bound already proved in this chapter,
with fixed denominator $10^{80}$, 120 logarithm-series terms, 24 initial
Hurwitz terms and 16 Bernoulli corrections. Its elementary symmetric
coefficient table is rebuilt through degree seven; this changes no
analytic hypothesis in the index-uniform remainder formula.

A second implementation independently extracts rising-factorial coefficients
and encloses all 126 signs. It uses denominator $10^{100}$, 128
logarithm-series terms, 32 initial Hurwitz terms and 16 Bernoulli corrections.
Its precision was increased because the delivered 55-digit default could
not resolve some new narrow brackets after very small powers were multiplied
by large factorial coefficients. Unresolved intervals were never accepted
as signs. Both replays use exact integer and fraction arithmetic throughout
acceptance. The complete endpoints, analytic error budgets, source digests
and descent signs are retained in
\texttt{verification/higher-zero-certificates.json} and
\texttt{verification/higher-zero-independent.json}.

The resulting first thresholds are
\[
 (K_1^{\rm cl},\ldots,K_7^{\rm cl})=(1,1,1,1,3,4,5).
\]
This finite sequence suggests further threshold questions; it does not
prove a formula in the index. In particular no claim that $K_n^{\rm cl}=n-2$
for every $n\ge5$ is made. Extending the terminating critical-point method
with effective saturation estimates is a concrete next research problem.
